% =====================================================================
%  chapters/minggu-02.tex
%  Minggu 2 — Figma Dasar & Pengantar UI/UX
% =====================================================================

\chapter{Figma Dasar \& Pengantar UI/UX}
\label{chap:minggu-02}

% ---------------------------------------------------------------------
% 1. CAPAIAN PEMBELAJARAN
% ---------------------------------------------------------------------
\begin{capaian}
Setelah menyelesaikan bab ini, mahasiswa diharapkan mampu:
\begin{itemize}
  \item mengoperasikan antarmuka dasar Figma (frame, shape, text, layers);
  \item membuat wireframe sederhana untuk satu halaman web;
  \item menjelaskan tiga konsep dasar UI/UX: heuristik, hierarchy, dan aksesibilitas;
  \item menerapkan konsep UI/UX ke dalam desain wireframe.
\end{itemize}
\end{capaian}

% ---------------------------------------------------------------------
% 2. TUJUAN PRAKTIK
% ---------------------------------------------------------------------
\begin{tujuanpraktik}
Pada sesi praktik ini, mahasiswa akan:
\begin{itemize}
  \item mengenal antarmuka dan alat dasar Figma;
  \item membuat wireframe halaman ``Profil Diri'' untuk perangkat \emph{mobile};
  \item mempelajari tiga konsep UI/UX dan menerapkannya pada wireframe.
\end{itemize}
\end{tujuanpraktik}

% ---------------------------------------------------------------------
% 3. LIVE CODING: PENGENALAN FIGMA
% ---------------------------------------------------------------------
\section{Pengenalan Figma}
\label{sec:m02-figma}

\subsection{Akses Figma}

\begin{langkah}
  \item Buka \texttt{https://www.figma.com} di peramban.
  \item Klik \textbf{Log in}, masuk dengan akun yang sudah dibuat (atau daftar jika belum).
  \item Kamu akan melihat halaman \textbf{Drafts} --- ini adalah ruang kerja \emph{default}.
\end{langkah}

\begin{tip}
Figma bisa digunakan langsung di peramban, tapi lebih nyaman menggunakan \textbf{Figma Desktop App} (unduh dari \texttt{https://www.figma.com/downloads/}).
\end{tip}

\subsection{Membuat File \& Frame Pertama}

\begin{langkah}
  \item Klik \textbf{New design file} (tombol biru di pojok kanan atas).
  \item Di \emph{canvas} kosong, tekan \textbf{F} di keyboard $\rightarrow$ arahkan mouse ke \emph{canvas} $\rightarrow$ \emph{drag} untuk membuat sebuah \textbf{Frame}.
  \item Di panel \textbf{Design} (kanan), atur ukuran frame: pilih preset \textbf{Web 1920} (1920\,$\times$\,1080\,px) atau \textbf{MacBook Air}.
\end{langkah}

\begin{tangkapanlayar}
  {assets/images/minggu-02/figma-frame-kosong.png}
  {Canvas Figma dengan satu frame kosong berukuran Web 1920.}
  {Pastikan panel Design di kanan menunjukkan ukuran frame yang benar.}
\end{tangkapanlayar}

\subsection{Alat Dasar Figma}

Pelajari toolbar di bagian atas \emph{canvas}:

\begin{table}[ht]
\centering
\caption{Shortcut alat dasar Figma}
\label{tab:m02-shortcut}
\begin{tabular}{llp{6cm}}
\toprule
\textbf{Shortcut} & \textbf{Alat} & \textbf{Fungsi} \\
\midrule
V & Move & Memilih \& memindahkan objek \\
F & Frame & Membuat \emph{frame} (artboard) \\
R & Rectangle & Membuat kotak persegi \\
O & Ellipse & Membuat lingkaran/oval \\
T & Text & Membuat teks \\
L & Line & Membuat garis \\
P & Pen & Membuat bentuk bebas \\
Ctrl + Z & Undo & Membatalkan aksi terakhir \\
\bottomrule
\end{tabular}
\end{table}

\textbf{Latihan cepat:} Di dalam \emph{frame} yang sudah dibuat, coba:
\begin{enumerate}
  \item Tekan \textbf{R} $\rightarrow$ \emph{drag} di dalam \emph{frame} untuk membuat kotak.
  \item Tekan \textbf{T} $\rightarrow$ klik di dalam \emph{frame} $\rightarrow$ ketik ``Hello Figma''.
  \item Tekan \textbf{V} $\rightarrow$ klik objek untuk memilih $\rightarrow$ \emph{drag} untuk memindahkan.
\end{enumerate}

\subsection{Panel Layers \& Properties}

Di panel kiri terdapat \textbf{Layers} --- setiap objek yang dibuat akan muncul sebagai \emph{layer}.

\begin{itemize}
  \item \textbf{Klik ganda} \emph{layer} untuk mengganti nama (misal: ``Rectangle 1'' menjadi ``Header Background'').
  \item \textbf{Urutan layer} = urutan di atas \emph{canvas} (layer di atas menimpa layer di bawah).
\end{itemize}

Di panel kanan terdapat \textbf{Design Properties} --- ubah:
\begin{itemize}
  \item \textbf{Fill:} warna isi objek
  \item \textbf{Stroke:} garis tepi
  \item \textbf{Corner Radius:} sudut membulat
  \item \textbf{Font \& Size:} untuk objek teks
\end{itemize}

\begin{tip}
Ubah warna kotak yang sudah dibuat menjadi biru, dan ubah teks ``Hello Figma'' menjadi ukuran 24\,px untuk berlatih mengatur \emph{properties}.
\end{tip}

% ---------------------------------------------------------------------
% 4. LIVE CODING: MEMBUAT WIREFRAME
% ---------------------------------------------------------------------
\section{Membuat Wireframe Halaman Web}
\label{sec:m02-wireframe}

\begin{catatan}
\textbf{Wireframe} adalah sketsa rendah detail yang menunjukkan struktur dan penempatan elemen halaman. Bukan desain final.
\end{catatan}

\begin{table}[ht]
\centering
\caption{Tingkatan detail desain}
\label{tab:m02-tingkatan}
\begin{tabular}{llp{6cm}}
\toprule
\textbf{Tingkat Detail} & \textbf{Nama} & \textbf{Fungsi} \\
\midrule
Rendah & Wireframe & Menunjukkan layout \& struktur (tanpa warna/font detail) \\
Sedang & Mockup & Menambahkan warna, font, gambar (visual lebih realistis) \\
Tinggi & Prototype & Interaktif --- bisa diklik dan dinavigasi \\
\bottomrule
\end{tabular}
\end{table}

Kita mulai dari \textbf{wireframe} (tingkat rendah) karena fokusnya adalah struktur, bukan tampilan cantik.

\subsection{Langkah 1: Buat Frame untuk \emph{Mobile}}

\begin{langkah}
  \item Tekan \textbf{F} $\rightarrow$ pilih preset \textbf{iPhone 14} (390\,$\times$\,844\,px) di panel kanan.
  \item \emph{Rename layer frame} ini menjadi \texttt{Wireframe - Profil Diri (Mobile)}.
\end{langkah}

\subsection{Langkah 2: Header / Navigasi}

\begin{langkah}
  \item Tekan \textbf{R} $\rightarrow$ buat \emph{rectangle} di bagian atas \emph{frame} (lebar penuh 390\,px, tinggi \~{}60\,px). \emph{Rename} layer jadi \texttt{Header}. Warna: abu-abu terang (\#E0E0E0).
  \item Tekan \textbf{T} $\rightarrow$ klik di atas \emph{rectangle} header $\rightarrow$ ketik \texttt{Nama Kamu}. Posisikan di kiri \emph{header}. Font: regular, 16\,px.
\end{langkah}

\subsection{Langkah 3: Hero / Foto Profil}

\begin{langkah}
  \item Tekan \textbf{O} $\rightarrow$ buat lingkaran di bawah \emph{header} (120\,$\times$\,120\,px, posisi center horizontal). Warna: abu-abu (\#CCCCCC). \emph{Rename} layer: \texttt{Foto Profil}.
  \item Tekan \textbf{T} $\rightarrow$ di bawah lingkaran, ketik nama lengkap dan keterangan program studi. Font: \textbf{bold} untuk nama, regular untuk keterangan.
\end{langkah}

\subsection{Langkah 4: Bagian ``Tentang Saya''}

\begin{langkah}
  \item Tekan \textbf{T} $\rightarrow$ ketik \texttt{Tentang Saya} sebagai judul \emph{section}. Font: bold, 18\,px.
  \item Tekan \textbf{R} $\rightarrow$ buat 2--3 \emph{rectangle} kecil di bawah judul sebagai \emph{placeholder} paragraf (lebar \~{}350\,px, tinggi \~{}16\,px $\times$ 3 baris). Warna: abu-abu sangat terang (\#F0F0F0).
\end{langkah}

\subsection{Langkah 5: Bagian ``Keahlian''}

\begin{langkah}
  \item Ketik judul \texttt{Keahlian}.
  \item Buat 3--4 \emph{rectangle} kecil sebagai \emph{placeholder} ``skill card'' (ukuran \~{}160\,$\times$\,80\,px, sejajar horizontal atau \emph{grid} 2$\times$2). Di dalam masing-masing, tambahkan teks \texttt{Skill 1}, \texttt{Skill 2}, dst.
\end{langkah}

\subsection{Langkah 6: Footer}

\begin{langkah}
  \item Tekan \textbf{R} $\rightarrow$ buat \emph{rectangle} di bagian bawah \emph{frame} (lebar penuh, tinggi \~{}50\,px). Warna: abu-abu gelap (\#999999).
  \item Tambahkan teks: \texttt{© 2026 Nama Kamu}.
\end{langkah}

\begin{tangkapanlayar}
  {assets/images/minggu-02/wireframe-profil-diri.png}
  {Wireframe halaman ``Profil Diri'' di Figma.}
  {Pastikan terlihat section: Header, Foto Profil, Tentang Saya, Keahlian, dan Footer dengan rapi.}
\end{tangkapanlayar}

% ---------------------------------------------------------------------
% 5. TIGA KONSEP DASAR UI/UX
% ---------------------------------------------------------------------
\section{Tiga Konsep Dasar UI/UX}
\label{sec:m02-uiux}

\subsection{Heuristik (Evaluasi Kegunaan)}

\textbf{Heuristik} adalah aturan umum desain yang membantu agar antarmuka pengguna (\emph{UI}) mudah digunakan. Berikut tiga heuristik paling relevan untuk pemula:

\begin{table}[ht]
\centering
\caption{Tiga heuristik utama}
\label{tab:m02-heuristik}
\small
\begin{tabular}{lp{4cm}p{5cm}}
\toprule
\textbf{Heuristik} & \textbf{Penjelasan} & \textbf{Contoh Penerapan di Wireframe} \\
\midrule
Visibility of System Status & Pengguna harus tahu apa yang sedang terjadi & Tambahkan indikator \emph{loading} saat halaman memuat \\
Match Between System and Real World & Gunakan bahasa yang dipahami pengguna & Judul ``Tentang Saya'' lebih jelas daripada ``Bio Data Seksi 4'' \\
Error Prevention & Cegah kesalahan pengguna sebelum terjadi & Form input harus punya \emph{placeholder}/label yang jelas \\
\bottomrule
\end{tabular}
\end{table}

\begin{catatan}
Ada 10 heuristik Nielsen Norman Group secara lengkap. Tiga di atas adalah yang paling sering ditemui dan paling mudah diterapkan untuk pemula.
\end{catatan}

\subsection{Hierarchy (Hierarki Visual)}

\textbf{Hierarchy} adalah pengaturan elemen agar mata pembaca tahu mana yang dilihat duluan.

\begin{table}[ht]
\centering
\caption{Jenis hierarki visual}
\label{tab:m02-hierarchy}
\begin{tabular}{llp{5cm}}
\toprule
\textbf{Jenis} & \textbf{Fungsi} & \textbf{Contoh} \\
\midrule
Hierarki ukuran & Elemen lebih besar = lebih penting & Nama besar di atas, deskripsi kecil di bawah \\
Hierarki warna & Warna kontras menarik perhatian & Tombol ``Hubungi Saya'' berwarna biru di antara elemen abu-abu \\
\bottomrule
\end{tabular}
\end{table}

\textbf{Latihan di Figma:} Pastikan di wireframe yang sudah dibuat:
\begin{enumerate}
  \item Nama di bagian Hero \textbf{lebih besar} dari teks deskripsi.
  \item Judul section (Tentang Saya, Keahlian) \textbf{lebih besar} dari konten di bawahnya.
  \item Jika belum, ubah ukuran teks di Figma untuk mencerminkan hierarchy.
\end{enumerate}

\subsection{Aksesibilitas (\emph{Accessibility})}

\textbf{Aksesibilitas} adalah desain yang bisa digunakan oleh semua orang, termasuk yang memiliki keterbatasan.

\begin{table}[ht]
\centering
\caption{Aspek aksesibilitas untuk pemula}
\label{tab:m02-akses}
\begin{tabular}{lp{8cm}}
\toprule
\textbf{Aspek} & \textbf{Yang Harus Diperhatikan} \\
\midrule
Kontras warna & Teks harus terbaca dengan latar belakang. Rasio kontras minimal \textbf{4.5:1} untuk teks normal. \\
Ukuran teks & Minimal \textbf{16\,px} untuk teks \emph{body}. Jangan gunakan ukuran yang terlalu kecil. \\
Struktur heading & Gunakan heading secara berurutan (\texttt{<h1>} $\rightarrow$ \texttt{<h2>} $\rightarrow$ \texttt{<h3>}), jangan loncat-loncat. \\
\bottomrule
\end{tabular}
\end{table}

\begin{tip}
Gunakan \textbf{WebAIM Contrast Checker} (\texttt{https://webaim.org/resources/contrastchecker/}) untuk mengecek kontras warna.
\end{tip}

\textbf{Latihan di Figma:}
\begin{enumerate}
  \item Di wireframe, pastikan warna teks di atas \emph{background} cukup kontras.
  \item Jika menggunakan warna abu-abu di atas putih, pastikan tidak terlalu tipis/pucat.
\end{enumerate}

% ---------------------------------------------------------------------
% 6. LATIHAN MANDIRI
% ---------------------------------------------------------------------
\section{Latihan Mandiri}
\label{sec:m02-latihan}

\begin{latihan}[Wireframe Halaman ``Portofolio'' (Desktop)]
Buat wireframe baru untuk halaman portofolio di Figma dengan spesifikasi:
\begin{itemize}
  \item \textbf{Frame:} Web 1920 (1920\,$\times$\,1080\,px).
  \item \textbf{Section yang wajib ada:}
  \begin{enumerate}
    \item \textbf{Navbar} --- logo di kiri, 3--4 menu di kanan (Home, Portofolio, About, Contact).
    \item \textbf{Hero} --- judul besar + tombol ``Lihat Proyek''.
    \item \textbf{Grid Proyek} --- minimal 6 kotak (\emph{card}) yang mewakili proyek.
    \item \textbf{Tentang Saya} --- foto \emph{placeholder} + teks deskripsi.
    \item \textbf{Footer} --- hak cipta \& ikon media sosial \emph{placeholder}.
  \end{enumerate}
\end{itemize}
\end{latihan}

\begin{tangkapanlayar}
  {assets/images/minggu-02/wireframe-portofolio.png}
  {Wireframe halaman Portfolio di Figma (tampilan penuh).}
  {Pastikan terlihat Navbar, Hero, Grid Proyek, Tentang Saya, dan Footer.}
\end{tangkapanlayar}

\begin{latihan}[Evaluasi Heuristik pada Website Nyata]
\begin{enumerate}
  \item Buka salah satu website berikut: \texttt{google.com}, \texttt{tokopedia.com}, atau \texttt{instagram.com} (halaman login).
  \item Evaluasi website tersebut menggunakan tiga heuristik yang sudah dipelajari dan catat hasilnya:
\end{enumerate}
\end{latihan}

\begin{table}[ht]
\centering
\caption{Template evaluasi heuristik}
\label{tab:m02-eval}
\begin{tabular}{p{4.5cm}cp{5cm}}
\toprule
\textbf{Heuristik} & \textbf{Ya/Tidak} & \textbf{Bukti / Contoh} \\
\midrule
Visibility of System Status & & \\
Match Between System and Real World & & \\
Error Prevention & & \\
\bottomrule
\end{tabular}
\end{table}

\begin{latihan}[Cek Kontras Aksesibilitas]
\begin{enumerate}
  \item Pilih 2 kombinasi warna dari wireframe yang kamu buat.
  \item Buka \texttt{https://webaim.org/resources/contrastchecker/}.
  \item Masukkan warna teks dan \emph{background}, cek rasio kontras.
  \item Pastikan minimal satu kombinasi memenuhi \textbf{WCAG AA} (4.5:1 untuk teks normal).
\end{enumerate}
\end{latihan}

\begin{tangkapanlayar}
  {assets/images/minggu-02/webaim-contrast-checker.png}
  {WebAIM Contrast Checker dengan hasil pengecekan.}
  {Pastikan terlihat rasio kontras dan status ``Pass'' atau ``Fail''.}
\end{tangkapanlayar}

% ---------------------------------------------------------------------
% 7. TROUBLESHOOTING
% ---------------------------------------------------------------------
\section{Troubleshooting}
\label{sec:m02-troubleshoot}

\begin{table}[ht]
\centering
\caption{Masalah umum Figma dan solusinya}
\label{tab:m02-troubleshoot}
\small
\begin{tabular}{p{4.5cm}p{8cm}}
\toprule
\textbf{Masalah} & \textbf{Solusi} \\
\midrule
Figma tidak bisa dibuka di peramban & Gunakan Figma Desktop App, atau coba peramban Chrome/Firefox versi terbaru \\
Frame tidak bisa dibuat & Tekan shortcut \textbf{F} (huruf F, bukan \emph{function key}) untuk aktifkan alat Frame \\
Objek tidak bisa dipindahkan & Pastikan alat \textbf{Move (V)} aktif, bukan alat lain \\
Layer tidak muncul di panel & \emph{Scroll} panel Layers di sebelah kiri \emph{canvas} --- mungkin tersembunyi di bawah \\
Perubahan tidak tersimpan & Figma \emph{auto-save}, tapi pastikan koneksi internet stabil. Cek ikon \emph{cloud} di pojok kanan atas \\
Tidak bisa membuat duplikat & Pilih objek $\rightarrow$ tekan \textbf{Ctrl + D} (Windows) atau \textbf{Cmd + D} (Mac) \\
\bottomrule
\end{tabular}
\end{table}

% ---------------------------------------------------------------------
% 8. RANGKUMAN
% ---------------------------------------------------------------------
\section{Rangkuman}
\label{sec:m02-rangkuman}

\begin{itemize}
  \item Mengoperasikan Figma: membuat \emph{frame}, \emph{shapes}, teks, mengatur \emph{properties}.
  \item Membuat wireframe sederhana untuk halaman web (5 section: Header, Foto, Tentang Saya, Keahlian, Footer).
  \item Tiga konsep dasar UI/UX: Heuristik, Hierarchy, Aksesibilitas.
  \item Menerapkan \emph{hierarchy} dan aksesibilitas di wireframe.
  \item Mengecek kontras warna dengan WebAIM Contrast Checker.
\end{itemize}
